\section{Introduction}

Financial markets represent a challenging domain for artificial intelligence research, characterized by high volatility, low signal-to-noise ratio, non-stationarity of raw data, and complex decision-making dynamics under uncertainty. Deep reinforcement learning has emerged as a promising approach for sequential decision-making problems, yet its application to portfolio allocation remains underexplored in modular architectures. This thesis does not aim to generate novel alpha or outperform established financial benchmarks. Rather, it investigates whether decomposing the portfolio allocation problem into specialized modules can improve learning stability and efficiency.

A key proposition motivates this work: passive buy-and-hold strategies for broad-market ETFs may be inferior to active strategies by ignoring daily price signals. By using end-of-day market data and computational models to dynamically adjust portfolio weights, investors can capture prevailing price trends and better manage downside risk \cite{Faber2007}, \cite{Moskowitz2012}. However, extracting stationary as well as memory-preserving  trading signals from non-stationary financial data while simultaneously determining optimal portfolio weights remains difficult. Traditional approaches, such as static asset allocation or purely predictive forecasting, often fail to adapt to regime changes or capture the inherently sequential nature of trading decisions.

The transition from static forecasting to sequential decision-making in non-stationary markets requires a framework that can adapt continuously. Deep reinforcement learning in combination with the right algorithms is well-suited to this task because it naturally models sequential decisions, handles continuous action spaces, and learns policies that generalize beyond training data. Yet practitioners face two fundamental challenges when applying DRL to portfolio management. First, the action space is inherently high-dimensional and continuous: the agent must specify target allocations across multiple assets simultaneously. Second, financial data exhibits an extremely low signal-to-noise ratio, where genuine economic signals are obscured by microstructure noise and random walk dynamics.

This thesis proposes addressing these challenges through hierarchical decomposition. Layer 1: Single-Asset Agent (SAA) extracts temporal features from individual assets using recurrent neural networks, learning asset-specific temporal dynamics in isolation. Layer 2: Portfolio Allocator Agent (PAA) uses these pre-extracted features alongside raw market data to determine optimal portfolio weights via self-attention mechanisms. By separating temporal extraction from cross-sectional allocation, each module can specialize to its respective task. This raises a central research question: \textit{Does modularizing temporal (single-asset) and cross-sectional (portfolio-level) functions improve the stability and efficiency of deep reinforcement learning agents for dynamic portfolio allocation?}

The proposed architecture is evaluated through walk-forward out-of-sample testing and multiple baseline validation comparisons to distinguish genuine alpha from memorized patterns. By isolating the contribution of each module, this work aims to provide empirical evidence for hierarchical decomposition as a viable paradigm for financial reinforcement learning.